% Figure from MATH 493 notes, section 31.
\begin{tikzpicture}[pt/.style={circle, draw, thick, minimum size=0.62cm, inner sep=0pt, font=\small},
  ar/.style={-{Stealth[length=2.2mm]}, thick}]
  % ---------- sigma ----------
  \begin{scope}
    \node[font=\small] at (1.1,2.05) {$\sigma = (1\,2\,3)(4\,5)$};
    \node[pt, figB] (a1) at (0,1.2) {$1$};
    \node[pt, figB] (a2) at (0.75,0) {$2$};
    \node[pt, figB] (a3) at (-0.75,0) {$3$};
    \node[pt, figB] (a4) at (2.1,1.2) {$4$};
    \node[pt, figB] (a5) at (2.1,0) {$5$};
    \draw[ar] (a1) to[bend left=18] (a2);
    \draw[ar] (a2) to[bend left=18] (a3);
    \draw[ar] (a3) to[bend left=18] (a1);
    \draw[ar] (a4) to[bend left=30] (a5);
    \draw[ar] (a5) to[bend left=30] (a4);
  \end{scope}
  % ---------- relabelling ----------
  \draw[-{Stealth[length=2.5mm]}, thick, black!60] (3.0,0.6) -- (4.6,0.6)
    node[midway, above, font=\small, black] {rename} node[midway, below, font=\small, black] {$i \mapsto \tau(i)$};
  % ---------- tau sigma tau^{-1} ----------
  \begin{scope}[shift={(6.3,0)}]
    \node[font=\small] at (1.1,2.05) {$\tau\sigma\tau^{-1} = (2\,5\,1)(3\,4)$};
    \node[pt, figR] (b1) at (0,1.2) {$2$};
    \node[pt, figR] (b2) at (0.75,0) {$5$};
    \node[pt, figR] (b3) at (-0.75,0) {$1$};
    \node[pt, figR] (b4) at (2.1,1.2) {$3$};
    \node[pt, figR] (b5) at (2.1,0) {$4$};
    \draw[ar] (b1) to[bend left=18] (b2);
    \draw[ar] (b2) to[bend left=18] (b3);
    \draw[ar] (b3) to[bend left=18] (b1);
    \draw[ar] (b4) to[bend left=30] (b5);
    \draw[ar] (b5) to[bend left=30] (b4);
    \node[font=\scriptsize, black!60] at (1.1,-0.75) {$\tau = (1\,2\,5\,4\,3)$};
  \end{scope}
\end{tikzpicture}
